\section{Introduction}
\label{sec:introduction}

Stochastic filtering concerns the estimation of a time-varying hidden
state from noisy observations. The problem appears in many settings where
the quantity of interest cannot be measured directly but must instead be
inferred sequentially as new data arrive.

For linear systems with Gaussian noise, the filtering distribution can be
computed exactly using the Kalman filter. More general models typically
require numerical approximations. Particle filtering provides one such
approach by representing the filtering distribution with a collection of
weighted samples, or particles. Increasing the number of particles
improves the Monte Carlo approximation, but also increases the
computational cost.

Particle filters are particularly interesting from a parallel-computing
perspective. Propagation and likelihood evaluation can be carried out
independently across particles, while operations such as normalization,
estimation, and resampling require some form of global coordination.
The algorithm therefore contains both naturally parallel work and
potential scalability bottlenecks.

This project studies that interaction between filtering algorithms and
parallel computation. The work is organized in stages, with each stage
building on the previous implementation.

Stage~I uses a scalar linear-Gaussian model for which the Kalman filter
provides an exact benchmark. A bootstrap particle filter is first
validated against this reference solution and is then parallelized using
OpenMP. The resulting implementation is used to study race conditions,
reductions, strong scaling, and the effect of serial components on
end-to-end performance.

Later stages will extend the filtering and parallel-computing framework
as the project develops. The common questions are how accurately the
filter approximates the desired state distribution, which parts of the
algorithm can be parallelized effectively, and which operations ultimately
limit scalability.

Section~\ref{sec:background} summarizes the filtering concepts needed
for Stage~I. The remaining sections describe the individual stages,
their implementations, and their numerical and performance results.